\hypertarget{python-3.8-walrus-operator-and-positional-only-parameters}{%
\chapter{Python 3.8: Walrus Operator (:=) and Positional-Only Parameters
(/)}\label{python-3.8-walrus-operator-and-positional-only-parameters}}

\hypertarget{the-assignment-expression-pep-572}{%
\subsection{\texorpdfstring{14.1 The Assignment Expression (\texttt{:=}
/ PEP
572)}{14.1 The Assignment Expression (:= / PEP 572)}}\label{the-assignment-expression-pep-572}}

Introduced in Python 3.8, the assignment expression operator
(colloquially called the \textbf{walrus operator}) allows assigning
values to variables inside expressions. This departs from Python's
historical strict division between statements and expressions.

\hypertarget{ast-representation}{%
\subsubsection{1. AST Representation}\label{ast-representation}}

A standard assignment is a statement represented by the
\passthrough{\lstinline!Assign!} node in the Abstract Syntax Tree (AST),
which does not return a value. The walrus operator compiles to a
\passthrough{\lstinline!NamedExpr!} node, which is an expression node
returning the bound value.

\begin{lstlisting}
      [Assign Statement]                  [NamedExpr Expression]
         /         \                            /         \
   Name(x)      Constant(1)               Name(x)      Constant(1)
   (Returns void/no value)              (Returns 1 to parent context)
\end{lstlisting}

\hypertarget{bytecode-stack-operations}{%
\subsubsection{2. Bytecode \& Stack
Operations}\label{bytecode-stack-operations}}

Let's examine how CPython processes a normal assignment vs.~an
assignment expression. Consider the following source codes and their
corresponding bytecodes:

\begin{lstlisting}[language=Python]
# Standard Assignment
x = 1
\end{lstlisting}

\begin{lstlisting}
2 LOAD_CONST               0 (1)
4 STORE_FAST               0 (x)
\end{lstlisting}

\begin{lstlisting}[language=Python]
# Assignment Expression
(x := 1)
\end{lstlisting}

\hypertarget{cpython-3.8-bytecode}{%
\paragraph{CPython 3.8 Bytecode:}\label{cpython-3.8-bytecode}}

\begin{lstlisting}
2 LOAD_CONST               0 (1)
4 DUP_TOP
6 STORE_FAST               0 (x)
\end{lstlisting}

\hypertarget{cpython-3.11-optimized-bytecode}{%
\paragraph{CPython 3.11+ Optimized
Bytecode:}\label{cpython-3.11-optimized-bytecode}}

\begin{lstlisting}
2 LOAD_CONST               0 (1)
4 COPY                     1
6 STORE_FAST               0 (x)
\end{lstlisting}

The execution steps for the walrus operator stack operation are: 1.
\textbf{\passthrough{\lstinline!LOAD\_CONST!}}: Pushes the reference to
the constant integer \passthrough{\lstinline!1!} onto the value stack.
2. \textbf{\passthrough{\lstinline!DUP\_TOP!} /
\passthrough{\lstinline!COPY 1!}}: Replicates the top element of the
stack, resulting in two references to \passthrough{\lstinline!1!} on the
stack. 3. \textbf{\passthrough{\lstinline!STORE\_FAST!}}: Pops one of
the references and binds it to the local variable name
\passthrough{\lstinline!x!}. 4. \textbf{Stack Residual}: The remaining
reference to \passthrough{\lstinline!1!} remains on the top of the
stack, allowing parent expressions (such as \passthrough{\lstinline!if!}
conditionals or loops) to consume it.

\hypertarget{scoping-rules-and-symbol-table-traversal}{%
\subsubsection{3. Scoping Rules and Symbol Table
Traversal}\label{scoping-rules-and-symbol-table-traversal}}

To prevent variable leakage and namespace pollution, PEP 572 specifies
complex scoping rules, especially within comprehensions: *
\textbf{Comprehensions}: Python list, set, and dict comprehensions, as
well as generator expressions, are executed in a nested function scope.
* \textbf{Variable Binding (Hoisting)}: An assignment expression
\passthrough{\lstinline!x := ...!} inside a comprehension binds the
variable \passthrough{\lstinline!x!} in the \emph{surrounding} (parent)
scope, not the local comprehension scope. * \textbf{Symbol Table
Resolution}: During compilation, the symbol table builder
(\passthrough{\lstinline!symtable.c!}) traverses variables inside
comprehensions. When it encounters a
\passthrough{\lstinline!NamedExpr!}, it marks the target variable as a
\textbf{free variable} in the comprehension scope and a \textbf{cell
variable} in the parent scope. This hoists the reference out of the
comprehension's inner namespace.

Let's examine this hoisting in action. Consider the following code:

\begin{lstlisting}[language=Python]
def parent_func(data):
    result = [y for x in data if (y := x * 2) > 2]
    return result, y
\end{lstlisting}

\hypertarget{disassembly-of-parent_func}{%
\paragraph{\texorpdfstring{Disassembly of
\texttt{parent\_func}:}{Disassembly of parent\_func:}}\label{disassembly-of-parent_func}}

\begin{lstlisting}
  2           0 LOAD_CLOSURE             0 (y)          /* Load the cell reference for y */
              2 BUILD_TUPLE              1
              4 LOAD_CONST               1 (<code object <listcomp>...>)
              6 LOAD_CONST               2 ('parent_func.<locals>.<listcomp>')
              8 MAKE_FUNCTION            8 (closure)    /* Bind cell to function closure */
             10 LOAD_FAST                0 (data)
             12 GET_ITER
             14 CALL_FUNCTION            1
             16 STORE_FAST               1 (result)

  3          18 LOAD_FAST                1 (result)
             20 LOAD_DEREF               0 (y)          /* Load y from closure cell */
             22 BUILD_TUPLE              2
             24 RETURN_VALUE
\end{lstlisting}

\hypertarget{disassembly-of-the-nested-listcomp-code-object}{%
\paragraph{\texorpdfstring{Disassembly of the nested
\texttt{\textless{}listcomp\textgreater{}} code
object:}{Disassembly of the nested \textless listcomp\textgreater{} code object:}}\label{disassembly-of-the-nested-listcomp-code-object}}

\begin{lstlisting}
  2           0 BUILD_LIST               0
              2 LOAD_FAST                0 (.0)         /* Load incoming iterator */
        >>    4 FOR_ITER                26 (to 32)
              6 STORE_FAST               1 (x)
              8 LOAD_FAST                1 (x)
             10 LOAD_CONST               1 (2)
             12 BINARY_MULTIPLY
             14 DUP_TOP
             16 STORE_DEREF              0 (y)          /* Store into outer scope cell */
             18 LOAD_CONST               2 (2)
             20 COMPARE_OP               4 (>)
             22 POP_JUMP_IF_FALSE        4
             24 LOAD_DEREF               0 (y)          /* Load y from outer cell */
             26 LIST_APPEND              2
             28 JUMP_ABSOLUTE            4
        >>   32 RETURN_VALUE
\end{lstlisting}

By using \passthrough{\lstinline!STORE\_DEREF!} and
\passthrough{\lstinline!LOAD\_DEREF!}, the nested list comprehension
writes directly to the parent scope cell, allowing
\passthrough{\lstinline!y!} to survive after the list comprehension has
finished executing.

\hypertarget{compiler-restrictions}{%
\subsubsection{4. Compiler Restrictions}\label{compiler-restrictions}}

To prevent syntactically ambiguous or unstable code, the compiler
enforces strict boundaries: * \textbf{Iteration Variable Conflicts}: A
target variable cannot share a name with a comprehension iteration
variable. E.g.
\passthrough{\lstinline![i for i in range(10) if (i := 2)]!} raises
\passthrough{\lstinline!SyntaxError: assignment expression cannot rebind comprehension iteration variable!}.
* \textbf{Class Scope Prohibitions}: The walrus operator is explicitly
blocked inside comprehensions in class bodies:
\passthrough{\lstinline!python     class MyClass:         data = [1, 2, 3]         result = [y for x in data if (y := x * 2) > 2]!}
This raises
\passthrough{\lstinline!SyntaxError: assignment expression within a comprehension cannot be used in a class body!}.
\emph{Why?} Class namespaces are evaluated as temporary dict namespaces,
not standard function frames. They do not support closure cells
(\passthrough{\lstinline!PyCellObject!}). Since the comprehension runs
as a separate nested helper function, it cannot bind closures to class
body locals, creating a scoping mismatch.

\begin{center}\rule{0.5\linewidth}{0.5pt}\end{center}

\hypertarget{positional-only-parameters-pep-570}{%
\subsection{\texorpdfstring{14.2 Positional-Only Parameters (\texttt{/}
/ PEP
570)}{14.2 Positional-Only Parameters (/ / PEP 570)}}\label{positional-only-parameters-pep-570}}

Python 3.8 introduced the \passthrough{\lstinline!/!} marker to define
positional-only parameters. Any parameters declared before the
\passthrough{\lstinline!/!} marker cannot be passed as keyword
arguments.

\hypertarget{c-level-representation-in-pycodeobject}{%
\subsubsection{\texorpdfstring{1. C-level Representation in
\texttt{PyCodeObject}}{1. C-level Representation in PyCodeObject}}\label{c-level-representation-in-pycodeobject}}

The compiler encodes parameter boundaries directly into the function's
bytecode descriptor struct. In \passthrough{\lstinline!code.h!}, the
\passthrough{\lstinline!PyCodeObject!} structure contains specific
integer fields to track positional boundaries:

\begin{lstlisting}[language=C]
typedef struct {
    PyObject_HEAD
    int co_argcount;            /* Number of positional and positional-only arguments */
    int co_posonlyargcount;     /* Number of positional-only arguments */
    int co_kwonlyargcount;      /* Number of keyword-only arguments */
    /* ... additional fields ... */
} PyCodeObject;
\end{lstlisting}

For a function definition:

\begin{lstlisting}[language=Python]
def func(a, b, /, c, d, *, e, f):
    pass
\end{lstlisting}

The compiler maps the arguments as follows: *
\passthrough{\lstinline!a, b!}: Positional-only.
(\passthrough{\lstinline!co\_posonlyargcount = 2!}) *
\passthrough{\lstinline!c, d!}: Positional-or-keyword. (Total positional
\passthrough{\lstinline!co\_argcount = 4!}, representing positional-only
+ positional-or-keyword) * \passthrough{\lstinline!e, f!}: Keyword-only.
(\passthrough{\lstinline!co\_kwonlyargcount = 2!})

\hypertarget{argument-parsing-and-vectorcall-optimization-pep-590}{%
\subsubsection{2. Argument Parsing and Vectorcall Optimization (PEP
590)}\label{argument-parsing-and-vectorcall-optimization-pep-590}}

When a function is called, CPython passes arguments using the
\textbf{Vectorcall} calling protocol introduced in Python 3.8:

\begin{lstlisting}[language=C]
PyObject *vectorcall(PyObject *callable, PyObject *const *args, size_t nargsf, PyObject *kwnames);
\end{lstlisting}

Where: * \passthrough{\lstinline!args!} is a pointer to a contiguous
array containing all passed argument values. *
\passthrough{\lstinline!nargsf!} specifies the number of positional
arguments. * \passthrough{\lstinline!kwnames!} is a tuple of strings
containing the keyword argument names. The keyword values are stored in
the \passthrough{\lstinline!args!} array starting at index
\passthrough{\lstinline!nargsf!}.

\hypertarget{the-vectorcall-matching-algorithm}{%
\paragraph{The Vectorcall Matching
Algorithm:}\label{the-vectorcall-matching-algorithm}}

\begin{enumerate}
\def\labelenumi{\arabic{enumi}.}
\tightlist
\item
  \textbf{Keyword Verification}: CPython parses the keyword strings in
  \passthrough{\lstinline!kwnames!}. If any string matches a parameter
  name located before index
  \passthrough{\lstinline!co\_posonlyargcount!} (e.g.~trying to pass
  \passthrough{\lstinline!a=1!}), the interpreter raises a
  \passthrough{\lstinline!TypeError!}.
\item
  \textbf{Bypassing Dictionary Checks}: Because positional-only
  parameters are guaranteed to be passed positionally, CPython
  completely bypasses string-matching hash loops for the first
  \passthrough{\lstinline!co\_posonlyargcount!} arguments.
\item
  \textbf{Direct Array Offset Writing}: The VM copies references
  directly from the \passthrough{\lstinline!args!} array into the local
  frame variable cells by array index offset:
  \[\text{LocalFrameOffset}[i] = args[i] \quad \text{for } 0 \le i < co\_posonlyargcount\]
  This direct offset copy cuts down calling overhead by \textbf{10\% to
  15\%}, which is highly beneficial for builtins and low-level utility
  functions that are called repeatedly.
\end{enumerate}

\begin{center}\rule{0.5\linewidth}{0.5pt}\end{center}
